% TOP_PROBLEM: {"id": 7200079, "title": "Does every higher-dimensional tiling by translations of convex polytope tiles have an affine transformation taking it to a Voronoi diagram", "category_id": 6, "category": {"id": 6, "name": "geometry", "display_name": "Geometry", "description": "Euclidean and non-Euclidean geometry, geometric structures.", "slug": "geometry", "order_index": 6, "created_at": "2026-07-31T15:26:25.671Z"}, "status": "partially_solved"}
% FIRST_POSTED: 2026-09-27
% !TeX program = lualatex
% ATTEMPT_STATUS: unresolved
\documentclass[11pt]{article}
\usepackage[margin=25mm]{geometry}
\usepackage{fontspec}
\setmainfont{Latin Modern Roman}
\usepackage{amsmath,amssymb,mathtools}
\usepackage{unicode-math}
\setmathfont{Latin Modern Math}
\usepackage{xurl,hyperref}
\hypersetup{colorlinks=true,urlcolor=blue}
\setcounter{secnumdepth}{0}
\setlength{\parindent}{0pt}
\setlength{\parskip}{5pt}
\setlength{\emergencystretch}{3em}
\begin{document}
\section*{\texorpdfstring{Does every higher-dimensional tiling by translations of convex polytope tiles have an affine transformation taking it to a Voronoi diagram}{Does every higher-dimensional tiling by translations of convex polytope tiles have an affine transformation taking it to a Voronoi diagram}}\label{7200079}
\subsection{Definitions and mathematical statement}
A parallelohedron is a full-dimensional convex polytope $P\subset{\mathbb{R}}^d$ that tiles ${\mathbb{R}}^d$ by translations, with disjoint interiors of distinct tiles. For a full-rank lattice $\Lambda$, its Euclidean Dirichlet--Voronoi cell at zero is
\[
 V(\Lambda)=\{x\in{\mathbb{R}}^d:\|x\|\le\|x-v\|
                 \text{ for every }v\in\Lambda\}.
\]
The conjecture asks whether every parallelohedron is of the form $A V(\Lambda)+b$ for some invertible linear map $A$, vector $b$, and lattice $\Lambda$. An equivalent viewpoint allows a positive-definite quadratic form in place of the Euclidean norm. The requested equivalence is affine, not necessarily an isometry or a similarity. Finding a translational tiling alone does not identify it as a Voronoi cell after a linear change of coordinates.
\par\smallskip\textit{Formulation and concept references:} \href{https://arxiv.org/abs/1906.05193}{[S1]}.

\subsection{Short English statement}
Can every convex polytope that tiles space by translation be transformed affinely into a lattice's nearest-point region?
\subsection{Research attempt}
\paragraph{Scope and literature check.}
This attempt by GPT-6 Astra Ultra was prepared on 27 September 2026.
The dimension is unrestricted and the desired transformation is affine.
Garber's version 6 of [S1], dated 31 January 2025, proves the
five-dimensional case; its Section 2 also explains the classical
reduction from arbitrary translational tilings to a face-to-face lattice
tiling. Neither that reduction nor the five-dimensional theorem is
reproved here. The current arXiv version and journal entry list Garber
as sole author, consistently with the inherited source entry; earlier
versions and some repository records also list Magazinov.

The route below is to turn an already identified tiling lattice and its
facet data into a positive-definite matrix problem. The local-to-global
step can be proved exactly using volume. We then solve that matrix
problem for parallelotopes and for every centrally symmetric convex
hexagon in the plane, including an explicit formula. These are known
affirmative families, rederived to test the method. The unresolved step
is existence of a compatible positive-definite matrix for arbitrary
parallelohedra in arbitrary dimension.

\paragraph{Lemma 477.1: changing the metric is exactly the affine freedom.}
For a symmetric positive-definite matrix $Q$, write
\[
 \|x\|_Q^2=x^{\mathsf T}Qx,\qquad
 V_Q(\Lambda)=\{x:2x^{\mathsf T}Qt\le t^{\mathsf T}Qt
                                      \text{ for all }t\in\Lambda\}.
\]
Choose an invertible matrix $L$ with $Q=L^{\mathsf T}L$. Expanding
$\|x-t\|_Q^2-\|x\|_Q^2$ shows
\begin{equation}\label{a477:affine}
                        L V_Q(\Lambda)=V(L\Lambda).
\end{equation}
Conversely, pulling the Euclidean norm back through an invertible linear
map gives precisely such a positive-definite form. Translation first
puts the central symmetry centre of the tile at zero. Thus constructing
$Q$ for the centred polytope is sufficient for the catalogue conclusion.

Positive definiteness cannot be replaced by invertibility or by positive
values on a few selected vectors. It makes the distance coercive and
ensures that nearest lattice points exist. An indefinite quadratic form
does not define a nearest-point cell in this sense.

\paragraph{A volume fact used in the local-to-global argument.}
For every full-rank lattice $\Lambda$ and positive-definite $Q$,
\begin{equation}\label{a477:volume}
                    \operatorname{vol}V_Q(\Lambda)=\det\Lambda,
\end{equation}
where $\det\Lambda$ denotes its positive covolume. To justify this,
coercivity and lattice discreteness give a nearest lattice point to
every $x$. Points with two nearest lattice points lie in a countable
union of affine bisector hyperplanes, hence a null set. The translates
of $V_Q(\Lambda)$ therefore cover and have disjoint interiors.
The cell is bounded: choose a bounded fundamental parallelepiped, move
$x$ into it by a lattice translate, and use its finite $Q$-diameter as
a uniform bound on distance to the lattice. If $x\in V_Q$, zero is
nearest, so this bound also bounds $\|x\|_Q$.

Integrating
$\sum_{t\in\Lambda}\mathbf1_{V_Q(\Lambda)}(x-t)=1$ almost everywhere
over a fundamental parallelepiped proves \eqref{a477:volume}. The same
argument shows $\operatorname{vol}P=\det\Lambda$ whenever $P+\Lambda$
is a translational tiling with disjoint interiors. This proof does not
assume the conjecture: it uses a metric Voronoi cell only after a
positive-definite $Q$ has been specified.

\paragraph{Lemma 477.2: a finite facet certificate is sufficient.}
Let $P$ be centred at zero and have a face-to-face lattice tiling
$P+\Lambda$. For each facet $F$, let $t_F\in\Lambda$ be the centre
of its neighbouring tile and let $a_F$ be an outward normal, so
\begin{equation}\label{a477:facets}
 P=\bigcap_F\left\{x:a_F\cdot x\le\tfrac12a_F\cdot t_F\right\}.
\end{equation}
The half-support value follows from central symmetry: the shared facet
$P\cap(P+t_F)$ is preserved by $x\mapsto t_F-x$, so its centre is
$t_F/2$. The support constant is positive because zero is in the
interior.

Suppose there are a symmetric positive-definite matrix $Q$ and positive
scalars $\lambda_F$ such that
\begin{equation}\label{a477:certificate}
                              Qt_F=\lambda_F a_F
                                  \qquad\text{for every facet }F.
\end{equation}
Then $P=V_Q(\Lambda)$.

\emph{Proof.} The lattice inequalities for the particular vectors $t_F$
are exactly the inequalities in \eqref{a477:facets}, after multiplication
by $\lambda_F>0$. Since a Voronoi cell satisfies the inequalities for
\emph{all} lattice vectors, $V_Q(\Lambda)\subseteq P$. Both sets have
volume $\det\Lambda$ by \eqref{a477:volume}. A proper inclusion of
compact full-dimensional convex sets has a positive-volume difference:
a separating hyperplane and a small interior ball produce a region of
positive volume in the larger set. Thus inclusion with equal volumes
forces equality. This proves the lemma.

This argument closes a possible infinite-constraint gap. It would be
incorrect simply to assume that checking neighbouring lattice vectors
automatically checks all lattice vectors for an arbitrary norm. Here
the missing implication follows from a specified positive-definite
quadratic form, inclusion in the correct direction, and equal volumes.

If $P$ is already $V_Q(\Lambda)$ with its Voronoi lattice tiling, the
shared facet with $P+t_F$ lies on its $Q$-bisector, so its outward
normal is a positive multiple of $Qt_F$. Thus the certificate is also
necessary for this fixed tiling. The classical tiling reduction cited
in [S1] supplies the canonical face-to-face lattice when the initial
tiling has offsets between rows or other non-face-to-face features.
Our argument does not silently assume that the particular translate set
initially handed to us is itself a lattice.

\paragraph{Attempt to solve the general matrix system.}
The equations in \eqref{a477:certificate} are linear in the entries of
$Q$ and the $\lambda_F$, supplemented by symmetry, positivity of the
scalars, and $Q\succ0$. Multiplying $Q$ and all $\lambda_F$ by the
same positive scalar changes no cell, so one may impose
$\operatorname{tr}Q=1$. For exact rational facet data, a rational
candidate can be checked using those linear equations and Sylvester's
criterion, namely positivity of all leading principal minors of the
symmetric matrix. That is a finite certificate for an individual tile.

Finding a nonzero solution to the linear equations is not enough.
For example, prescribe in the plane
\[
 t_1=(1,0),\quad t_2=(0,1),\qquad
 a_1=(1,2),\quad a_2=(2,1).
\]
Both $a_i\cdot t_i$ are positive. If $Qt_i=\lambda_i a_i$ with
$\lambda_i>0$, symmetry forces $\lambda_1=\lambda_2=\lambda$ and
\[
                   Q=\lambda\begin{pmatrix}1&2\\2&1\end{pmatrix}.
\]
Its determinant is $-3\lambda^2$, so it is indefinite. These data do
not arise from the asserted facet neighbours of a tile for
$\mathbb Z^2$: the associated parallelogram
\[
          |x+2y|\le1/2,\qquad|2x+y|\le1/2
\]
has area $1/3$, not covolume $1$. It is itself a parallelogram and
does tile by a different lattice. Thus this is a counterexample to a
local linear-algebra shortcut, not to Voronoi's conjecture. It identifies
where actual tiling compatibility must enter any general proof.

Numerical feasibility has a related limitation. Eigenvalues that appear
positive at floating point precision are not an exact certificate near
the boundary of the positive cone. Conversely, failure of a search
restricted to diagonal matrices or a prescribed symmetry class does not
rule out an unrestricted $Q$. The conjecture requires an existence
argument, not the success of a particular search parameterization.

\paragraph{Proposition 477.3: every parallelotope has an explicit metric.}
Let $B$ be any invertible real $d\times d$ matrix and put
\[
 P=B[-1/2,1/2]^d,\qquad\Lambda=B\mathbb Z^d.
\]
For any positive diagonal matrix $D$, take
\begin{equation}\label{a477:boxmetric}
                             Q=B^{-\mathsf T}DB^{-1}.
\end{equation}
For $x=Bu$ and $t=Bn$, one has
$\|x-t\|_Q^2=\sum_jD_{jj}(u_j-n_j)^2$. Each term is minimized
independently by a nearest integer. Zero is a simultaneous minimizer
exactly when $|u_j|\le1/2$ for every $j$. Therefore
$V_Q(\Lambda)=P$, including boundary ties.

For a sheared square, $B=\left(\begin{smallmatrix}1&s\\0&1\end{smallmatrix}\right)$
and $D=I$ give
\[
              Q=\begin{pmatrix}1&-s\\-s&1+s^2\end{pmatrix},
              \qquad\det Q=1.
\]
Its off-diagonal entries are part of the solution. Requiring the
Euclidean metric for the original lattice, or requiring $Q$ to be
diagonal in the original coordinates, would unnecessarily discard
these valid examples.

\paragraph{Proposition 477.4: an explicit metric for every central hexagon.}
Let $P$ be a strictly convex centrally symmetric hexagon, with vertices
in counterclockwise order
$a,b,c,-a,-b,-c$. Here strict convexity means that consecutive edges
are not collinear; a degenerate six-vertex presentation reduces to a
parallelogram. The centre is in the interior. The vectors $a,c$ are
independent and their positive cone contains $b$. After an invertible
linear change of coordinates, write
\[
                a=(1,0),\qquad c=(0,1),\qquad b=(u,v).
\]
Counterclockwise turning at $a,b,c$, respectively, gives
\begin{equation}\label{a477:hexconditions}
       u>0,\quad v>0,\quad u+v>1,\quad |u-v|<1.
\end{equation}
For example, the three edge cross products are $1+v-u$, $u+v-1$,
and $1+u-v$. Conversely, these inequalities give the indicated convex
hexagon with zero in its interior.

Define
\begin{equation}\label{a477:hexmetric}
 h=\frac{1-u^2-v^2}{2uv},\qquad
 Q=\begin{pmatrix}1&h\\h&1\end{pmatrix}.
\end{equation}
The inequalities \eqref{a477:hexconditions} are exactly what is needed
for $-1<h<1$: these are equivalent to
$(u-v)^2<1<(u+v)^2$. Hence $Q\succ0$. Furthermore
\[
                    a^{\mathsf T}Qa=b^{\mathsf T}Qb
                                           =c^{\mathsf T}Qc=1.
\]
The six vertices lie on a centred ellipse for this metric.

Set
\[
        t_1=a+b,\qquad t_3=c-a,\qquad
        \Lambda=\mathbb Zt_1+\mathbb Zt_3,\qquad t_2=b+c=t_1+t_3.
\]
The three positive-edge lines are the $Q$-bisectors of $0$ and
$t_1,t_2,t_3$, respectively. For the edge from $a$ to $b$,
\[
 (b-a)^{\mathsf T}Q(a+b)=b^{\mathsf T}Qb-a^{\mathsf T}Qa=0,
 \quad
 a^{\mathsf T}Q(a+b)=\tfrac12(a+b)^{\mathsf T}Q(a+b).
\]
The edge from $b$ to $c$ is identical with the letters shifted. For
the edge from $c$ to $-a$, use the sum $c-a=t_3$ and the equality
of the $Q$-norms of $c$ and $-a$. The opposite edges correspond to
$-t_1,-t_2,-t_3$. Since zero is in the interior and every right-hand
side is positive, the halfspaces have the required orientation.
It follows that $V_Q(\Lambda)\subseteq P$.

There is no assumption here that this candidate lattice already tiles
by $P$. Instead, calculate both volumes. The shoelace formula gives
\[
 \operatorname{area}P
    =\det(a,b)+\det(b,c)+\det(a,c)=u+v+1,
\]
while
\[
 \det\Lambda=\det(t_1,t_3)
     =\det\begin{pmatrix}1+u&-1\\v&1\end{pmatrix}=u+v+1.
\]
Equation \eqref{a477:volume} now forces $V_Q(\Lambda)=P$.
This simultaneously constructs the lattice tiling and proves the
affine Voronoi property. There is no circular use of a presumed tiling
by the lattice just constructed.

For $u=v=1$, the metric is
$Q=\left(\begin{smallmatrix}1&-1/2\\-1/2&1\end{smallmatrix}\right)$,
with determinant $3/4$, and the lattice basis is $(2,1),(-1,1)$.
The common area is $3$. For general $u,v$, the smallest eigenvalue is
$1-|h|$, so it can tend to zero as the hexagon approaches a degenerate
boundary. Strict inequalities are needed for each nondegenerate tile;
a semidefinite limit alone is not a certificate for it.

Combining this result with Proposition 477.3 covers the planar
classification of parallelohedra, namely parallelograms and centrally
symmetric hexagons, as recalled in [S1]. The classification is an
attributed input. The metrics and the volume proof for those two
families are supplied explicitly above.

\paragraph{A closure operation giving higher-dimensional examples.}
If $P_i=V_{Q_i}(\Lambda_i)$ for $i=1,2$, then
\begin{equation}\label{a477:product}
 P_1\times P_2
    =V_{\operatorname{diag}(Q_1,Q_2)}(\Lambda_1\times\Lambda_2).
\end{equation}
The squared distance is a sum of two independent squared distances.
If each coordinate is nearest to zero, every lattice pair has at least
as large a total distance. Conversely, testing lattice pairs with one
coordinate zero enforces the two separate conditions. Thus the equality
holds, and affine images of these products are also covered by
Lemma 477.1. Products of the explicit planar examples and intervals
give verified families in every dimension.

This does not exhaust higher-dimensional parallelohedra. In particular,
a proof that every tile decomposes into these products has not been
given and is not part of the cited tiling structure theorem. The
five-dimensional theorem in [S1] uses substantially richer local
combinatorics; its dimension-specific analysis is not replaced by
the elementary product construction.

\paragraph{Verification and the exact general gap.}
The rational checks accompanying this attempt test the explicit hexagon
formula on many positive rational $u,v$ satisfying
\eqref{a477:hexconditions}: positive determinant, equal vertex norms,
all three edge orthogonality relations, and equality of polygon area
with lattice covolume. These are checks of algebraic identities, not
a classification of tiles. The preceding symbolic proof establishes
the family for all admissible real $u,v$.

The proposed general route has two distinct tasks. First obtain the
facet-neighbour data of the canonical lattice tiling, using the known
structure theory. Then prove that its linear compatibility equations
have a solution meeting the positive-definite cone and all positive
facet scalings. Lemma 477.2 would finish the argument once that second
task were accomplished. Local positivity, isolated low-dimensional
sections, and examples from products do not establish it for arbitrary
facet systems. The indefinite-matrix diagnostic shows why positivity
cannot simply be omitted from the task.

\textbf{UNRESOLVED.} We have a proved finite certificate for a specified
tile and lattice, explicit solutions for parallelotopes and all planar
central hexagons, and closure under products and affine maps. Existence
of such a certificate for every parallelohedron in unrestricted
dimension is not proved here. The known five-dimensional theorem is
credited to [S1], and no new general solution is claimed.

\subsection{Sources}
\begin{itemize}
\item[S1]A. Garber, Voronoi conjecture for five-dimensional parallelohedra, Inventiones Mathematicae (2025); arXiv:1906.05193.
\url{https://arxiv.org/abs/1906.05193}.
\textit{Formulation source.}
\item[E1]Alexey Garber, \textit{Voronoi conjecture for five-dimensional
parallelohedra}, arXiv:1906.05193v6, 31 January 2025, Sections 1--2 and
Theorem 4.1; full-text version of [S1].
\url{https://arxiv.org/html/1906.05193v6}.
Used for the scope of the known theorem and the face-to-face lattice
reduction; the elementary matrix and planar arguments here are rederived.
\end{itemize}
\end{document}
